\documentclass[10pt]{article}
\usepackage[T1]{fontenc}
\renewcommand{\rmdefault}{ptm}
\usepackage[margin=0.75in]{geometry}
\usepackage{amsmath,amssymb,booktabs,tabularx,array,multirow,xcolor,float,hyperref}
\usepackage[expansion=false]{microtype}
\hypersetup{hidelinks,pdftitle={PRIME-Soft: Supplementary Design and Replay Details}}
\definecolor{TraceGray}{HTML}{929BA4}
\newcolumntype{Y}{>{\raggedright\arraybackslash}X}
\newcommand{\papertablestyle}{\fontsize{9}{11}\selectfont\setlength{\tabcolsep}{4pt}\renewcommand{\arraystretch}{1.05}}
\title{PRIME-Soft: Supplementary Design and Replay Details}
\author{}\date{}
\begin{document}
\maketitle
\section{Design-Space Domains and Parameter Bindings}

\label{app:design-space}
Table~\ref{tab:design-space} specifies the implemented base-only serial family through its generator domains. The family definition composes the base segment, guide, connector, and payload prototypes with the serial assembly rule. A study selects its permitted subset and freezes the associated task and control configuration.
\begin{table}[H]
\centering
\caption{Parameterized design space of the tendon-driven soft-robot family: implemented base-only serial configuration.}
\label{tab:design-space}
\papertablestyle
\begin{tabularx}{\textwidth}{@{}>{\itshape}>{\raggedright\arraybackslash}p{0.14\textwidth}!{\color{TraceGray}\vrule width 0.25pt}>{\raggedright\arraybackslash}p{0.24\textwidth}>{\raggedright\arraybackslash}p{0.10\textwidth}Y@{}}
\toprule
\normalfont\textbf{Category} & \textbf{Parameter} & \textbf{Type} & \textbf{Domain or assembly rule} \\
\midrule
Topology & Segment count; tendon count & Discrete & 1--9 segments; 1--7 tendons; consecutive base-origin spans \\
Geometry & Segment length; section family and dimensions & Mixed & 60--240~mm generic length; circle, tube, ellipse, rectangle; indexed bounds in the text below \\
Material & Young's modulus & Continuous & 2--15~MPa per segment \\
Tendon routing & Radius; guide and anchor station; angular placement & Mixed & Radius 14--23~mm; guide $s\in[0.15,0.75]$; anchor $s\in[0.8,1]$; phase $[-2\pi,2\pi]$; offsets $[-\pi,\pi]$~rad \\
Assembly & Tendon diameter; force rating; component placement & Mixed & Diameter 0.5--1.2~mm; force limit 1--8~N; guides and connectors at serial boundaries, payload at the tip \\
\midrule
Control $(U,\theta)$ & Input sequence; initialization; controller settings & Continuous & Task-specific admissible input; pretension 0--1~N; terminal/holding speed weights 0--1 in the base configuration \\
\bottomrule
\end{tabularx}
\end{table}


Indexed segment domains override the generic domain: segment~0 has length 80--240~mm, segment~1 has length 60--200~mm, and subsequent segments use 60--240~mm. Circular radius and elliptical semiaxes range from 3 to 12~mm; tube outer and inner radii range from 5 to 12~mm and 1 to 4~mm, respectively; rectangular width and height range from 6 to 18~mm. The compiler enforces positive section dimensions, tube wall thickness, ordered route stations, nonzero spans, owned guide holes, routing clearances, and consistent force and pretension limits. The base-only rule starts each tendon span at the base and routes it through consecutive segments to its anchor.

Density, bending viscosity, connector and payload inertia, and prototype station angles are inherited from the selected assembly in this configuration. Active pretension and controller weights belong to the control configuration; passive fixed-payout conditions are specified by the passive assembly and load model. Numerical cell count controls discretization. Parameter identifiers, units, and effective model bindings are resolved for each compiled candidate.

Throughout the sensitivity formulas, derivatives with respect to $d$ act only on $d_{\mathrm{cont}}$, with $d_{\mathrm{disc}}$ and the local parameterization held fixed. Likewise, design perturbation directions have zero discrete components.



\section{Recorded Model, Prompt, and Operation Configuration}

The archived baseline configuration and full-framework runtime identify the model as \texttt{gpt-5.6-sol} and configure medium reasoning. The parameter-tradeoff context replay uses the same model identifier with reasoning set to \texttt{none}; its fixed system prompt requests one next research decision and permits \texttt{research\_plan}, \texttt{research\_recall}, and \texttt{evidence\_read} within four provider requests.

{\raggedright The archived baseline command interface exposes \texttt{generate}, \texttt{evaluate}, and \texttt{finish} to both methods. Direct Mathematical Tools additionally exposes \texttt{math} with \texttt{jacobian}, \texttt{mechanics}, \texttt{finite\_difference}, \texttt{equilibrium}, and \texttt{kkt} operations. Its direct-optimization records additionally use nonlinear least squares for configuration fitting and bounded SLSQP for tendon-force allocation. The archived full-framework configuration permits structure generation, control initialization, physical replay, diagnosis, selection, and stopping, together with candidate-bound mechanism analysis, condition-wise control optimization, and comparison of saved control-adapted responses. Task configurations specify the enabled subset.\par}

The \href{https://github.com/wran021004-gif/softrobot-agent/tree/codex/gated-harness-paper-20261124/paper_handoff/supplementary_data/reproducibility}{supplementary reproducibility records} preserve the verbatim system messages from the archived baseline requests, the full-framework study instructions, and the context-replay system prompt and tool schemas, with source paths and file hashes.

\section{Context Evidence Checklist and Retrieval Records}
For parameter tradeoff (checkpoint PHS\_C037\_A), the five factual groups in Table~\ref{tab:context-evidence-checklist} were frozen before blind scoring. A group is covered only when the decision correctly uses its substantive content. Exact numbers and every listed source reference are not required; a partial reference counts if it preserves the relevant physical interpretation, whereas a bare citation or an incorrect interpretation does not. Coverage is binary per group and reported as covered groups divided by five, separately from the five-dimension decision score. Flat access covers G1--G5; hierarchical access covers G2, G4, and G5.

\begin{table}[H]
\centering
\caption{Predefined decision-critical evidence at the parameter-tradeoff checkpoint. Source identifiers refer to the frozen research archive.}
\label{tab:context-evidence-checklist}
\papertablestyle
\begin{tabularx}{\textwidth}{@{}>{\raggedright\arraybackslash}p{0.045\textwidth}>{\raggedright\arraybackslash}p{0.46\textwidth}Y@{}}
\toprule
\textbf{ID} & \textbf{Decision-critical evidence} & \textbf{Coverage criterion} \\
\midrule
G1 & Matched near-load results for viscosity variants F2/F3/F4 show reduced drift but nonmonotone speed and static-balance residuals (C029, C032, C035). & Recognizes the mixed trends rather than inferring uniform improvement from reduced drift. \\
G2 & Loaded qualification remains unsuccessful despite some unloaded or local mathematical qualification (C031, C034, C036, C037). & Distinguishes unloaded/local qualification from loaded dynamic qualification. \\
G3 & F3 has much larger finite-load motion transfer under distal loading than near loading (C032, C037). & Uses the near/distal contrast with the correct load-case and candidate scope. \\
G4 & Raw route records detect a cable--solid intersection at the initial pose/time (C031, C034, C036). & Separates installed-route invalidity from contact counts and late-time settling. \\
G5 & Sampled exploratory evidence covers selected cases, without establishing every-step qualification, full independent acceptance, numerical convergence, or a unique failure cause (C027, C031, C034, C036, C037). & Keeps the proposed decision within the demonstrated scope of the sampled evidence. \\
\bottomrule
\end{tabularx}
\end{table}

{\raggedright The hierarchical session makes three \texttt{research\_recall} calls with \texttt{view=details}: step~31 reads \texttt{C031\_EP0\_qualifications}, step~34 reads \texttt{C034\_EP0\_qualifications}, and step~36 reads \texttt{C036\_EP0\_route\_geometry}. Each call selects its record through the exact \texttt{/information/} pointer. The first two expose loaded-qualification records for the matched viscosity variants; the third exposes the installed-route geometry check. Coverage is assessed from the final decision, separately from whether a record was retrieved.\par}

\end{document}
